\begin{abstract}
Concept bottleneck models promise intrinsic interpretability: route the prediction through a
human-readable concept layer, and the concept activation \emph{is} the explanation. We build one
for dense prediction---a routing bottleneck that assigns every encoder token to one of $K$
clinically labeled experts---and train it for 3D brain-tumor segmentation. By every measure this
literature reports, it works: routing tracks the clinical concepts at $\casfg$ up to $0.96$ and
separates necrosis, edema and enhancing tumor at AUROC $0.88$, against $0.74$ for the strongest
post-hoc attribution. Yet the routing is \emph{causally inert}. Forcing tokens through a different
expert changes $1.1\%$ of voxel labels, and zeroing the bottleneck outright costs at most $3.1$
Dice---the decoder reconstructs the segmentation without ever consulting the concept layer.
Alignment metrics cannot detect this, because alignment is compatible with a bottleneck the model
never consults, and the alignment loss trains precisely what those metrics score. The cause is not
degenerate experts, which compute genuinely distinct transformations, but two leaks that
encoder--decoder architectures invite: a residual path \emph{through} the module and skip
connections \emph{around} it. We distill the diagnosis into three cheap tests---causal
intervention, bottleneck ablation, and a probe control---that expose the failure without
retraining, and then repair it. \emph{Removing} the decoder's alternatives does not work; the
optimizer compensates. Making the routing an \emph{additive term of the output} does: causal
control rises from $0.011$ to $\mathbf{0.636}$ $[0.621,0.651]$ label-flip under intervention
($p{<}10^{-28}$), the bottleneck turns from optional into indispensable (necrosis Dice
$0.48\!\to\!0.01$ without it), and the price is ${\approx}1$ Dice point. \textbf{A concept layer
is an explanation only if the model consults it, and showing that requires intervening, not
correlating.}
\end{abstract}
